\subsubsection*{Solution}
\paragraph*{Step-by-Step Derivation}

The first-harmonic continuum model and parameters used here agree with those employed in the tMoTe$_2$ literature; the corresponding plane-wave Hamiltonian is also given explicitly in the literature. \href{\detokenize{https://pmc.ncbi.nlm.nih.gov/articles/PMC10895274/}}{PNAS continuum-model formulation}, \href{\detokenize{https://www.nature.com/articles/s42005-025-02483-6}}{parameter cross-check}.

\subparagraph*{1. Moiré scales}

For $\theta=3.5^\circ$,

\[
a_M=\frac{3.52\ {\rm \text{\AA}}} {2\sin(1.75^\circ)} =57.6321442995\ {\rm \text{\AA}}.
\]

Consequently,

\[
b\equiv |\boldsymbol b_1| =\frac{4\pi}{\sqrt{3}a_M} =0.125888036011\ {\rm \text{\AA}}^{-1}.
\]

The reciprocal vectors can be written as

\[
\boldsymbol b_1=b(1,0),\qquad \boldsymbol b_2=b\left(\frac12,\frac{\sqrt3}{2}\right),
\]

and the Brillouin-zone area is

\[
A_{\rm BZ} = |\boldsymbol b_1\times\boldsymbol b_2| = \frac{\sqrt3}{2}b^2.
\]

The kinetic coefficient is

\[
\frac{\hbar^2}{2m^*} = \frac{1}{0.6}\frac{\hbar^2}{2m_e} = 6349.9701916\ {\rm meV\,\text{\AA}^2}.
\]

\subparagraph*{2. Plane-wave Hamiltonian}

Introduce the basis states $|l,\boldsymbol Q\rangle$, where

\[
\boldsymbol Q_t=\boldsymbol q_1+n\boldsymbol b_1+m\boldsymbol b_2,
\]

and

\[
\boldsymbol Q_b=-\boldsymbol q_1+n\boldsymbol b_1+m\boldsymbol b_2.
\]

Imposing

\[
|\boldsymbol Q_l|<4.1|\boldsymbol b_1|
\]

gives $63$ momenta in each layer and hence a $126\times126$ Hamiltonian at every $\boldsymbol k$.

With $s_b=-1$ and $s_t=+1$, the intralayer matrix elements are

\[
\adjustbox{max width=\linewidth}{$\displaystyle \begin{aligned}
H^{(l)}_{\boldsymbol Q,\boldsymbol Q'}(\boldsymbol k)
={}&
-\frac{\hbar^2}{2m^*}
|\boldsymbol k-\boldsymbol Q|^2
\delta_{\boldsymbol Q,\boldsymbol Q'}
\\
&+
V\sum_{i=1}^3
\left[
e^{is_l\psi}\,
\delta_{\boldsymbol Q',\,\boldsymbol Q+\boldsymbol g_i}
+
e^{-is_l\psi}\,
\delta_{\boldsymbol Q',\,\boldsymbol Q-\boldsymbol g_i}
\right].
\end{aligned}$}
\]

The interlayer block follows from the Fourier factor $e^{-i\boldsymbol q_i\cdot\boldsymbol r}$:

\[
H^{bt}_{\boldsymbol Q_b,\boldsymbol Q_t} = w\sum_{i=1}^3 \delta_{\boldsymbol Q_b-\boldsymbol Q_t,\boldsymbol q_i}, \qquad H^{tb}=(H^{bt})^\dagger.
\]

Diagonalizing this matrix at every momentum and sorting the energies from highest to lowest defines the top three bands.

\subparagraph*{3. Chern numbers}

For each isolated band, define the normalized lattice link

\[
U_\mu^{(n)}(\boldsymbol k) = \frac{ \langle u_{n,\boldsymbol k} | u_{n,\boldsymbol k+\Delta\boldsymbol k_\mu}\rangle }{ \left| \langle u_{n,\boldsymbol k} | u_{n,\boldsymbol k+\Delta\boldsymbol k_\mu}\rangle \right| },
\]

where

\[
\Delta\boldsymbol k_1=\frac{\boldsymbol b_1}{L}, \qquad \Delta\boldsymbol k_2=\frac{\boldsymbol b_2}{L}.
\]

The gauge-invariant plaquette flux is

\[
F_{12}^{(n)}(\boldsymbol k) = \operatorname{Arg} \left[ U_1^{(n)}(\boldsymbol k) U_2^{(n)}(\boldsymbol k+\Delta\boldsymbol k_1) U_1^{(n)}(\boldsymbol k+\Delta\boldsymbol k_2)^{-1} U_2^{(n)}(\boldsymbol k)^{-1} \right],
\]

and

\[
C_n=\frac{1}{2\pi}\sum_{\boldsymbol k}F_{12}^{(n)}(\boldsymbol k).
\]

At Brillouin-zone boundaries, the plane-wave coefficients are translated according to

\[
u_{l,\boldsymbol Q}(\boldsymbol k+\boldsymbol b_\alpha) = u_{l,\boldsymbol Q-\boldsymbol b_\alpha}(\boldsymbol k).
\]

The resulting Chern numbers, from the top band downward, are

\[
(C_1,C_2,C_3)=(1,-1,0).
\]

Their sum vanishes,

\[
C_1+C_2+C_3=0,
\]

as expected for this three-band group. Reversing the valley or the Berry-curvature orientation gives the globally reversed sequence $(-1, 1, 0)$, consistent with reported tMoTe$_2$ conventions. \href{\detokenize{https://journals.aps.org/prx/pdf/10.1103/PhysRevX.13.031037}}{PRX topology convention}.

\subparagraph*{4. Integrated quantum metric of the top band}

Only the kinetic term depends explicitly on $\boldsymbol k$. With

\[
\kappa\equiv\frac{\hbar^2}{2m^*},
\]

its derivatives in the plane-wave basis are diagonal:

\[
D_i\equiv\frac{\partial H}{\partial k_i}
=-2\kappa(k_i-Q_i)\delta_{ll'}\delta_{\boldsymbol Q\boldsymbol Q'}.
\]

For an isolated band $n$, first-order perturbation theory gives the gauge-invariant metric

\[
g_{ij}^{(n)}(\boldsymbol k)
=\operatorname{Re}\sum_{m\ne n}
\frac{\langle u_n|D_i|u_m\rangle\langle u_m|D_j|u_n\rangle}
{(E_n-E_m)^2}.
\]

Thus the trace needed for the top band is evaluated directly from the full eigensystem at each mesh point:

\[
\operatorname{Tr}g^{(n)}(\boldsymbol k)
=\sum_{m\ne n}
\frac{|\langle u_m|D_x|u_n\rangle|^2+|\langle u_m|D_y|u_n\rangle|^2}
{(E_n-E_m)^2}.
\]

This avoids the finite-link bias of estimating derivatives from neighboring eigenvector overlaps. The requested $60\times60$ Brillouin-zone quadrature is

\[
\operatorname{Tr}\mathcal G
\simeq \frac{A_{\rm BZ}}{60^2}
\sum_{l_1,l_2}\operatorname{Tr}g(\boldsymbol k_{l_1l_2}).
\]

Numerically,

\[
\operatorname{Tr}\mathcal G=7.51.
\]

This quantity is dimensionless because $g_{ij}$ has units of length squared while $d^2k$ has units of inverse length squared.
